% Generated by scripts/assemble.py; corrections belong in source/corrections.json.
\examhead{187}{陕西师范大学 2024 年数学分析试卷}{176}
% p177-b19; source PDF page 177
\blockanchor{p177-b19}
\question{1}(15 分) 判断对错并说明理由(正确的给出证明, 错误的给出反例).\par
% p177-b20; source PDF page 177
\blockanchor{p177-b20}
(1) 若 $f(x)$ 在 $(a, b)$ 上连续, 则 $f(x)$ 在 $(a, b)$ 上有界.\par
% p177-b21; source PDF page 177
\blockanchor{p177-b21}
(2) 若 $f(x)$ 在 $(a, b)$ 上连续, 且 $\lim_{x \to a^{+}} f(x)$ 与 $\lim_{x \to b^{-}} f(x)$ 均存在, 则 $f(x)$ 在 $(a, b)$ 上有界.\par
% p177-b22; source PDF page 177
\blockanchor{p177-b22}
(3) 若 $f(x)$ 在 $(a, b)$ 上连续且有界, 则 $\lim_{x \to a^{+}} f(x)$ 与 $\lim_{x \to b^{-}} f(x)$ 均存在.\par
% p177-b23; source PDF page 177
\par\Needspace{6\baselineskip}
\blockanchor{p177-b23}
\question{2}(15 分) 数列\par
\nopagebreak[4]
% p177-b24; source PDF page 177
\blockanchor{p177-b24}
\displayformula{a _ {1} = 2, a _ {2} = 2 + \frac {1}{2}, a _ {3} = 2 + \frac {1}{2 + \frac {1}{2}}, a _ {4} = 2 + \frac {1}{2 + \frac {1}{2 + \frac {1}{2}}}, \dots .}
% p177-b25; source PDF page 177
\blockanchor{p177-b25}
问数列 $\{a_{n}\}$ 是否收敛？说明理由.\par
% p177-b26; source PDF page 177
\par\Needspace{6\baselineskip}
\blockanchor{p177-b26}
\question{3}(15 分) 已知\par
\nopagebreak[4]
% p177-b27; source PDF page 177
\blockanchor{p177-b27}
\displayformula{f (x) = \left\{ \begin{array}{l l} (x - 1) ^ {3} \cos \frac {1}{x}, & x <   1; \\ a, & x = 1; \\ e ^ {b (x - 1)} + c, & x > 1. \end{array} \right.}
% p177-b28; source PDF page 177
\blockanchor{p177-b28}
问： $a,b,c$ 为何值时， $f(x)$ 在 $x = 1$ 处可导，并计算 $f^{\prime}(1)$\par
% p178-b00; source PDF page 178
\blockanchor{p178-b00}
\question{4}(15 分) 已知 $z = z(x, y)$ 是由方程 $x + 2y + 3z = e^z$ 确定的隐函数, 计算 $\frac{\partial z}{\partial x}, \frac{\partial z}{\partial y}, \frac{\partial^2 z}{\partial x^2}, \frac{\partial^2 z}{\partial y \partial x}$.\par
% p178-b01; source PDF page 178
\blockanchor{p178-b01}
\question{5}(15 分) 求抛物线 $y^{2} = px, y^{2} = qx$ ($0 < p < q$) 及双曲线 $xy = a, xy = b$ ($0 < a < b$) 所围成区域 $D$ 的面积.\par
% p178-b02; source PDF page 178
\par\Needspace{6\baselineskip}
\blockanchor{p178-b02}
\question{6}(15 分) 计算\par
\nopagebreak[4]
% p178-b03; source PDF page 178
\blockanchor{p178-b03}
\displayformula{\iint_ {\Sigma} \frac {(x + y) \mathrm{d} y \mathrm{d} z + y \mathrm{d} z \mathrm{d} x + 2 z \mathrm{d} x \mathrm{d} y}{\sqrt {x ^ {2} + y ^ {2} + z ^ {2}}}.}
% p178-b04; source PDF page 178
\blockanchor{p178-b04}
其中 $\Sigma$ 为球面 $x^{2} + y^{2} + z^{2} = a^{2}(a > 0)$ 的外侧.\par
% p178-b05; source PDF page 178
\blockanchor{p178-b05}
\question{7}(15 分) 设 $f$ 在闭区间 $[a, b]$ 上有定义, 且对任意的 $x, y \in [a, b]$, 有 $|f(x) - f(y)| \leq 2(x - y)^2$, 证明: $f(x)$ 在 $[a, b]$ 上为常值函数.\par
% p178-b06; source PDF page 178
\blockanchor{p178-b06}
\question{8}(15 分) 求幂级数 $\sum_{n=1}^{\infty} \frac{\ln(n+1)}{n} x^n$ 的收敛域.\par
% p178-b07; source PDF page 178
\par\Needspace{6\baselineskip}
\blockanchor{p178-b07}
\question{9}(15 分) 设三元函数 $f(x, y, z)$ 在 $D = \{(x, y, z) \mid x + y + z \leq 1\}$ 上可微, 且 $\forall (x, y, z) \in D$, 有\par
\nopagebreak[4]
% p178-b08; source PDF page 178
\blockanchor{p178-b08}
\displayformula{\left| f _ {x} (x, y, z) \right| \leq 1, \left| f _ {y} (x, y, z) \right| \leq 2, \left| f _ {z} (x, y, z) \right| \leq 3.}
% p178-b09; source PDF page 178
\par\Needspace{6\baselineskip}
\blockanchor{p178-b09}
证明： $\forall (x_{1},y_{1},z_{1}),(x_{2},y_{2},z_{2})\in D,$ 有\par
\nopagebreak[4]
% p178-b10; source PDF page 178
\blockanchor{p178-b10}
\displayformula{\left| f \left(x _ {2}, y _ {2}, z _ {2}\right) - f \left(x _ {1}, y _ {1}, z _ {1}\right) \right| \leq \left| x _ {2} - x _ {1} \right| + 2 \left| y _ {2} - y _ {1} \right| + 3 \left| z _ {2} - z _ {1} \right|.}
% p178-b11; source PDF page 178
\blockanchor{p178-b11}
\question{10}(15 分) 设 $f(x)$ 在闭区间 $[a, b]$ 上只有可去间断点, 证明: $f(x)$ 在 $[a, b]$ 上有界.\par
